\documentclass[conference]{IEEEtran}
\IEEEoverridecommandlockouts

\usepackage[T1]{fontenc}
\usepackage[utf8]{inputenc}
\usepackage[brazil]{babel}
\usepackage{cite}
\usepackage{amsmath,amssymb,amsfonts}
\usepackage{graphicx}
\usepackage{textcomp}
\usepackage{xcolor}
\usepackage{booktabs}
\usepackage{url}
\usepackage{array}

\newcommand{\todo}[1]{\textcolor{red}{[PREENCHER: #1]}}
\newcommand{\sensecv}{\textit{SenseCV}}
\renewcommand{\abstractname}{Resumo}
\renewcommand{\IEEEkeywordsname}{Palavras-chave}
\renewcommand{\refname}{Referências}

\def\BibTeX{{\rm B\kern-.05em{\sc i\kern-.025em b}\kern-.08em
    T\kern-.1667em\lower.7ex\hbox{E}\kern-.125emX}}

\begin{document}

\title{SenseCV: Uma Base de Dados Visual--Inercial Egocêntrica e Multilocal para Navegação Assistiva em Ambientes Internos}

\author{ \IEEEauthorblockN{ André Luis Albuquerque Pinheiro, Isaque Efraim Martins Oliveira,\\Elias Teodoro da Silva Jr. } \IEEEauthorblockA{ Computer Science Department - Federal Institute of Ceará (IFCE) Fortaleza, CE, Brazil\\ \{andre.luis.albuquerque07, isaque.efraim94\}@aluno.ifce.edu.br,  elias@ifce.edu.br } }

\maketitle

\begin{abstract}

\end{abstract}

\begin{IEEEkeywords}

\end{IEEEkeywords}

\section{Introdução}

A deficiência visual pode comprometer a autonomia em atividades de orientação e locomoção, especialmente em ambientes desconhecidos, dinâmicos ou sujeitos a alterações. A Organização Mundial da Saúde estima que ao menos 2,2 bilhões de pessoas apresentam algum grau de deficiência visual para perto ou para longe \cite{who2026}. Para essa população, obstáculos no nível do solo ou suspensos, corredores estreitos, circulação de pessoas e mudanças na organização do ambiente podem tornar o deslocamento interno uma tarefa complexa e aumentar a dependência de auxílio externo.

Avanços em visão computacional e inteligência artificial têm impulsionado o desenvolvimento de tecnologias assistivas capazes de reconhecer caminhos livres, detectar obstáculos, estimar a posição do usuário e produzir instruções de navegação \cite{casanova2025, khan2022}. Entretanto, muitas soluções dependem de câmeras RGB-D, LiDAR, sensores ultrassônicos, dispositivos vestíveis personalizados ou unidades externas de processamento. Embora esses recursos forneçam informações detalhadas sobre o ambiente, eles também podem aumentar o custo, o peso e a complexidade de implantação do sistema \cite{gao2025}.

Os smartphones constituem uma alternativa atraente por reunirem câmera, sensores inerciais, capacidade de processamento e armazenamento em um único dispositivo amplamente disponível. Além disso, uma solução pode utilizar os sensores durante a construção da base e depender somente da câmera na etapa de inferência. Essa estratégia reduz os requisitos para uma futura implantação e permite aproveitar os sinais inerciais como fonte de supervisão. O AI Guide Dog, por exemplo, emprega vídeos egocêntricos de smartphone para prever direções viáveis de navegação, utilizando os sensores móveis durante a coleta e a geração dos rótulos \cite{jadhav2025}. De forma relacionada, \cite{medeiros2021} utilizaram imagens, acelerômetro e giroscópio de smartphones para modelar movimentos frontais e laterais associados ao desvio de obstáculos em ambientes internos.

Apesar desses resultados, a construção e a avaliação das bases continuam sendo desafios centrais. Gravações egocêntricas apresentam elevada redundância temporal, pois quadros consecutivos frequentemente exibem os mesmos corredores, obstáculos e pontos de vista. Em bases constituídas por quadros extraídos de vídeos, a divisão aleatória pode distribuir imagens com alta correlação espaço-temporal entre os subconjuntos de treinamento, validação e teste, causando vazamento de informação e superestimando o desempenho dos modelos \cite{figueiredo2024leakage}. Estudos sobre sistemas de alerta de colisão também mostram que diferenças de qualidade visual, origem ou composição das classes podem ser exploradas como atalhos pelos modelos, sem que o fenômeno de interesse seja efetivamente aprendido \cite{chitraranjan2025}. Além disso, bases obtidas a partir de caminhadas naturais tendem a apresentar predominância da ação frontal, enquanto ações importantes para a segurança, como desviar e parar, possuem menor representação \cite{jadhav2025}.

Este artigo apresenta o \sensecv{}, uma base visual--inercial egocêntrica, coletada em múltiplos ambientes internos e desenvolvida para apoiar pesquisas em navegação local e desvio de obstáculos. A proposta combina um aplicativo próprio para aquisição sincronizada de vídeo e sensores, um protocolo controlado e centrado em eventos e uma estratégia de particionamento que preserva a independência entre as cenas utilizadas no treinamento, na validação e no teste. As principais contribuições deste trabalho são:

\begin{itemize}
    \item um aplicativo Android de código aberto que registra, de forma sincronizada, vídeo, metadados do dispositivo e diferentes dados de sensores internos;

    \item uma base coletada em três instituições, contendo eventos controlados de caminho livre, desvio lateral e parada diante de obstáculos;

    \item um esquema de rotulagem que combina as anotações do protocolo com os sinais do acelerômetro e do giroscópio;

    \item uma estratégia de particionamento por gravação, cena e local, evitando que quadros temporalmente correlacionados sejam distribuídos entre diferentes subconjuntos; e

    \item um protocolo experimental para comparar a predição direta da ação com uma formulação multitarefa orientada à presença de obstáculos.
\end{itemize}

O restante deste artigo está organizado da seguinte forma. A Seção II apresenta os trabalhos relacionados sobre navegação assistiva, visão egocêntrica e bases visual--inerciais. A Seção III descreve a metodologia, incluindo o aplicativo SenseCV, o protocolo de coleta, o esquema de rotulagem, a formulação do problema, as estatísticas da base, o particionamento e as métricas de avaliação. A Seção IV apresenta os resultados experimentais. Por fim, a Seção V reúne as conclusões e indica as próximas etapas da pesquisa.


\section{Trabalhos Relacionados}

A navegação visual em ambientes internos abrange diferentes problemas, como localização, planejamento de rotas, reconhecimento de caminhos livres, detecção de obstáculos e geração de instruções. Revisões da área mostram que soluções baseadas em câmeras apresentam vantagens relacionadas ao custo e à disponibilidade dos dispositivos, mas ainda enfrentam desafios de iluminação, oclusão, custo computacional, posicionamento da câmera e generalização entre ambientes \cite{khan2022}. Nesse contexto, os trabalhos mais próximos ao \sensecv{} podem ser organizados em soluções móveis para percepção de obstáculos, métodos de predição egocêntrica, bases visual--inerciais e conjuntos de dados originalmente desenvolvidos para plataformas robóticas.

\cite{breve2020} propuseram uma solução orientada a smartphones que classifica imagens como caminho livre ou não livre, demonstrando a viabilidade de utilizar transferência de aprendizado em uma base relativamente pequena, composta por cenas internas e externas. Em uma formulação baseada na identificação explícita dos objetos, \cite{bastidas2025detectra} apresentaram a base Detectra, formada por 7.600 imagens e 76 classes de obstáculos coletadas em uma instituição de ensino para pessoas com deficiência visual. Um modelo YOLOv5s foi integrado a uma aplicação móvel com retorno auditivo e háptico. Entretanto, as imagens representam objetos isolados ou composições de objetos e não sequências egocêntricas associadas à ação de navegação que deve ser executada pelo usuário.

Soluções vestíveis podem incorporar informações mais completas sobre o ambiente, embora exijam hardware adicional. \cite{gao2025} desenvolveram o WOAD, que combina câmera infravermelha, sensor de profundidade, processamento em FPGA, smartphone e retorno auditivo e tátil. O sistema estima características como distância, direção e movimento dos obstáculos, apresentando elevada capacidade de desvio em experimentos internos e externos. Essa abordagem evidencia os benefícios da percepção multimodal, mas também contrasta com propostas baseadas exclusivamente em sensores presentes em smartphones comuns.

A predição egocêntrica da direção de caminhada é particularmente próxima ao objetivo do \sensecv{}. \cite{jadhav2025} formularam a navegação como uma classificação multirrótulo de \textsc{frente}, \textsc{esquerda} e \textsc{direita} a partir de vídeos capturados por um smartphone posicionado no tórax. Os sensores móveis são utilizados na geração dos rótulos, enquanto a inferência final depende apenas das imagens. A base possui aproximadamente 57 horas e foi particionada de modo a evitar sobreposição de cenas. Interseções internas, nas quais mais de uma direção era possível, receberam correções manuais. Apesar da diversidade da base, o trabalho não contempla uma classe de parada e reconhece dificuldades em cenas dinâmicas, caminhos bloqueados e estruturas ambíguas.

\cite{medeiros2021} também utilizaram imagens de smartphone e sinais inerciais para estudar o desvio de obstáculos em ambientes internos. A tarefa foi organizada em duas classificações, com cinco níveis de velocidade frontal e cinco níveis de velocidade lateral, avaliando arquiteturas como VGG16, ResNet50 e DroNet. O trabalho demonstra a possibilidade de empregar o movimento medido pelos sensores como supervisão das imagens. Entretanto, a divisão dos dados é descrita no nível das imagens, sem explicitar o agrupamento por gravação, percurso ou ambiente. O \sensecv{} mantém a aquisição visual--inercial, mas utiliza ações locais discretas e preserva as gravações completas durante o particionamento.

\cite{plikynas2022} apresentaram uma proposta de navegação interna baseada em rotas coletadas por voluntários, câmeras RGB e de profundidade, IMU, SLAM, processamento em nuvem e retorno auditivo e tátil. Um aspecto relevante é o emparelhamento automático entre imagens e classes de movimento extraídas da IMU. Os autores utilizaram uma ConvLSTM para classificar os movimentos \textsc{frente}, \textsc{esquerda} e \textsc{direita}, além de estimar a confiabilidade da predição. Embora essa estratégia seja próxima à rotulagem visual--inercial, o sistema depende de equipamento vestível especializado, rotas previamente percorridas e infraestrutura computacional adicional. O \sensecv{}, por sua vez, concentra-se na disponibilização de dados capturados com smartphones comuns e em decisões locais que não dependem de uma rota previamente mapeada.

O aprendizado multitarefa fornece outra formulação relevante. \cite{loquercio2018} propuseram a DroNet, uma rede com extrator visual compartilhado e duas saídas: um valor contínuo de direção e uma probabilidade de colisão. \cite{chang2022} adaptaram essa formulação à base HDIN, coletada com VANTs em ambientes internos reais, e discutiram problemas como oscilações dos sensores, unidades incompatíveis e ruído nos rótulos contínuos. Embora as plataformas e os comandos de controle sejam distintos daqueles utilizados por pedestres, esses trabalhos motivam a comparação entre a predição direta de uma ação e o aprendizado conjunto da ação com uma representação auxiliar de obstáculo.

Bases capturadas por smartphones também são relevantes para tarefas diferentes da recomendação de direção. \cite{samadzadeh2024} apresentaram a AUT-VI, composta por 126 sequências capturadas em 17 locais, contendo vídeo, IMU e GPS sincronizados. A base inclui objetos dinâmicos, reflexos, mudanças de iluminação, movimentos bruscos e diferentes formas de segurar o aparelho, além de disponibilizar um aplicativo para reprodução do processo de coleta. Seu objetivo principal, entretanto, é a avaliação de algoritmos de odometria visual--inercial e estimação de pose, e não a escolha de uma ação local diante de obstáculos.

Conjuntos robóticos oferecem maior diversidade sensorial e de interações, mas não representam diretamente a dinâmica de uma pessoa com deficiência visual. \cite{karnan2022scand} apresentaram o SCAND, com 8,7 horas e 138 trajetórias coletadas por meio dos robôs Boston Dynamics Spot e Clearpath Jackal em situações sociais internas e externas. A base contém câmeras, LiDAR, IMU, odometria e comandos de joystick, além de interações com pessoas, multidões e veículos. Os robôs podem atingir velocidades de 1,6 e 2,0~m/s, respectivamente. Como referência, \cite{liu2019travel} consideram uma velocidade média de aproximadamente 0,72~m/s para a caminhada de pessoas com deficiência visual. Além da diferença de velocidade, a altura da câmera, a morfologia, a dinâmica e os comandos dos robôs limitam a aplicação direta do SCAND ao contexto assistivo humano.

\cite{yan2020flobot} apresentaram uma base coletada por um robô autônomo de limpeza em aeroporto, armazém e supermercado. O FLOBOT utiliza LiDAR 2D e 3D, câmeras RGB-D e câmera estéreo para detectar pessoas, objetos no chão e alterações do ambiente. Os experimentos mostram que modelos treinados em um local podem perder desempenho quando avaliados em outro, reforçando a importância da diversidade ambiental. Apesar dessa contribuição, a plataforma industrial, a quantidade de sensores e o objetivo de limpeza autônoma diferem da perspectiva egocêntrica e dos recursos disponíveis em um smartphone.

A Tabela~\ref{tab:related_comparison} resume os trabalhos mais próximos em relação à plataforma, ao ambiente, à tarefa e ao particionamento.

\begin{table*}[!t]
\centering
\caption{Comparação entre bases para navegação egocêntrica e assistiva.}
\label{tab:related_comparison}
\renewcommand{\arraystretch}{1.15}
\setlength{\tabcolsep}{5pt}
\scriptsize
\begin{tabular}{
p{2.4cm}
p{0.8cm}
p{2.4cm}
p{1.8cm}
p{5.1cm}
p{2.4cm}}
\hline
\textbf{Trabalho} &
\textbf{Ano} &
\textbf{Plataforma} &
\textbf{Ambiente} &
\textbf{Tarefa} &
\textbf{Particionamento} \\
\hline

Breve e Fischer \cite{breve2020} &
2020 &
Smartphone &
Interno e externo &
Classificação de caminho livre ou obstruído &
Não informado \\

Medeiros \textit{et al.} \cite{medeiros2021} &
2021 &
Smartphone &
Interno &
Classificação de velocidade frontal e lateral &
Por imagens \\

AIGD \cite{jadhav2025} &
2025 &
Smartphone &
Interno e externo &
Predição das direções frente, esquerda e direita &
Cenas disjuntas \\

Plikynas \textit{et al.} \cite{plikynas2022} &
2022 &
Dispositivo vestível &
Interno &
Predição das direções frente, esquerda e direita &
Por rotas \\

AUT-VI \cite{samadzadeh2024} &
2024 &
Smartphone &
Interno e externo &
Odometria visual--inercial &
Por sequências \\

Detectra \cite{bastidas2025detectra} &
2025 &
Smartphone &
Interno e externo &
Detecção de classes de obstáculos &
Por imagens \\

\textbf{SenseCV} &
\textbf{2026} &
\textbf{Smartphone} &
\textbf{Interno e multilocal} &
\textbf{Presença de obstáculo e ação local} &
\textbf{Por gravações e locais} \\

\hline
\end{tabular}
\end{table*}

Conforme sintetizado na Tabela~\ref{tab:related_comparison}, os trabalhos existentes apresentam contribuições importantes para detecção de obstáculos, predição egocêntrica, localização e aquisição visual--inercial. Entretanto, permanecem limitadas as bases que combinam aquisição por smartphone comum, vídeo egocêntrico sincronizado com sensores inerciais, múltiplos ambientes internos, eventos controlados de desvio e parada e particionamento sem sobreposição de cenas. O \sensecv{} busca preencher essa lacuna ao disponibilizar os dados brutos do conjunto de teste, o aplicativo de aquisição e um protocolo reproduzível para estudar decisões locais de navegação e a capacidade de generalização dos modelos.



\section{Metodologia}

A metodologia foi definida em torno de dois objetivos. Primeiro, a base deve expor os modelos a diferentes edifícios, corredores, pontos de vista, objetos e condições de iluminação. Segundo, a organização temporal dos dados deve ser preservada tanto na rotulagem quanto no particionamento. Os modelos baseados em imagens isoladas podem ser treinados e avaliados quadro a quadro, enquanto os modelos temporais utilizam sequências de quadros. Entretanto, em ambos os casos, todos os quadros pertencentes a uma mesma gravação ou evento devem permanecer no mesmo subconjunto, evitando que imagens consecutivas ou quase idênticas sejam distribuídas entre treino, validação e teste.

\subsection{Aplicativo de Coleta SenseCV}

A aquisição dos dados foi realizada com o aplicativo Android de código aberto \sensecv.\footnote{Repositório público contendo o código-fonte, as instruções de instalação e a documentação de uso do aplicativo: \url{https://github.com/SenseCV}.} Para cada gravação, o aplicativo armazena o vídeo da câmera, os timestamps associados aos quadros, os metadados do dispositivo e da configuração da câmera, além de arquivos no formato JavaScript Object Notation (JSON) contendo os dados dos sensores disponibilizados pelo aparelho. As modalidades registradas incluem acelerômetro, giroscópio e magnetômetro. Quando o serviço de localização está disponível, o aplicativo também coleta as coordenadas de latitude e longitude, seus respectivos timestamps e informações sobre o estado do GPS.

É importante distinguir as modalidades \emph{capturadas} daquelas \emph{utilizadas atualmente}. Embora o aplicativo armazene diferentes fluxos sensoriais para ampliar as possibilidades de reutilização da base, somente os dados do acelerômetro e do giroscópio são empregados no esquema de rotulagem adotado neste trabalho. Os dados de GPS e do magnetômetro são mantidos como informações complementares para estudos futuros, como análise de trajetórias, localização e fusão multimodal. Cada sessão é organizada em uma estrutura própria de arquivos no dispositivo, o que facilita a revisão, a seleção e, quando necessário, a anonimização das gravações antes de sua transferência ou disponibilização.

O arquivo de metadados registra informações como fabricante e modelo do dispositivo, versão do Android, câmera selecionada, distância focal, orientação, resolução, taxa de quadros e configurações da câmera. O usuário pode selecionar a lente, a orientação em modo retrato ou paisagem, o modo quadrado, a resolução e uma taxa nominal de 15, 30 ou 60~FPS. A gravação também pode ser iniciada e interrompida por meio de um controle Bluetooth externo, reduzindo a necessidade de interação direta com a tela e a movimentação provocada pelo toque no aparelho.

\subsection{Protocolo de Coleta}

As gravações foram realizadas em Fortaleza, Brasil, no Instituto Federal do Ceará (IFCE), na Universidade Federal do Ceará (UFC) e na Biblioteca Pública Estadual do Ceará (BECE). A versão atual contém 13 sessões coletadas entre junho e julho de 2026. A maioria das sessões utilizou um Samsung Galaxy A56, com lente traseira de 1,74~mm e vídeo quadrado de $480\times480$; uma sessão inicial utilizou um POCO X3 Pro em $720\times720$.

Em testes preliminares, observou-se que a fixação direta do smartphone ao corpo introduzia oscilações excessivas decorrentes da caminhada, afetando tanto a estabilidade do vídeo quanto os sinais inerciais. Para reduzir essas interferências e padronizar a orientação da câmera, as coletas foram realizadas com o\textbf{ajustar aqui, deixar mais claro} gimbal S5 CamfreeGo. O estabilizador diminui movimentos involuntários associados ao balanço corporal e permite distinguir com maior clareza as variações relacionadas às ações de seguir em frente, desviar e parar, aspecto particularmente relevante porque o acelerômetro e o giroscópio são empregados no esquema de rotulagem da base. \textbf{ajustar isso aqui} Em parte das coletas, foram utilizadas duas alturas nominais, aproximadamente 50~cm e 100~cm em relação ao piso, para introduzir variação no ponto de vista. \todo{confirmar o número total de portadores e descrever como o gimbal foi transportado e mantido nas alturas definidas}

O protocolo é centrado em eventos completos de navegação, e não em amostragem oportunista de quadros. Nos eventos de desvio, o portador inicia aproximadamente 5~m antes do obstáculo, começa a manobra a cerca de 1,5~m dele e continua por aproximadamente 1,5~m após ultrapassá-lo. Alguns corredores da BECE exigiram distâncias iniciais menores, entre 3,2 e 3,7~m, devido às limitações físicas do local. Nos eventos de parada, o portador se aproxima a partir de aproximadamente 5~m e para cerca de 1~m antes do obstáculo. Os eventos de caminho livre consistem em caminhada contínua sem obstáculo planejado, geralmente por dez ou mais passos.

O protocolo registra a posição do obstáculo, o desvio solicitado, o corredor, a orientação do celular e uma descrição textual do evento. Foram incluídos objetos estáticos comuns, como lixeiras, cones, placas de sinalização, caixas, cadeiras, extintores, escadas e livros, além de pessoas em diferentes poses. O desenho experimental produz deliberadamente repetições de desvios para a esquerda e para a direita, bem como sequências de parada e caminhada livre, reduzindo a predominância extrema da ação frontal encontrada em bases egocêntricas coletadas durante caminhadas naturais \cite{jadhav2025}.

Os dados brutos correspondentes exclusivamente ao conjunto de teste estão disponíveis publicamente no Kaggle.\footnote{Página do conjunto de dados, contendo as gravações originais e os arquivos sensoriais associados ao conjunto de teste: \url{https://www.kaggle.com/datasets/sensecv/test-set}.} Os conjuntos de treinamento e validação serão disponibilizados em trabalhos futuros\todo{é importante falar isso??}. A disponibilização inclui os vídeos, os dados dos sensores e os metadados gerados pelo aplicativo, preservando a organização original das sessões.

\subsection{Esquema de Rotulagem e Formulação do Problema}

Cada evento fornece dois alvos relacionados. O primeiro é um rótulo binário de obstáculo,

\begin{equation}
    o \in \{\textsc{livre},\textsc{obstáculo}\},
\end{equation}

em que os eventos de caminho livre recebem \textsc{livre}, enquanto os eventos com obstáculo à esquerda, ao centro, à direita ou com parada recebem \textsc{obstáculo}. O segundo é um rótulo de ação local com quatro classes,

\begin{equation}
    y \in \{\textsc{frente},\textsc{esquerda},\textsc{direita},\textsc{parar}\}.
\end{equation}

A ação segue a anotação do protocolo: ausência de desvio em caminho livre corresponde a \textsc{frente}; um desvio para a esquerda ou para a direita corresponde à respectiva ação; e um evento de aproximação encerrado antes de um obstáculo bloqueando o caminho corresponde a \textsc{parar}. Essa distinção impede que o lado do obstáculo seja confundido com a recomendação: um obstáculo à direita, por exemplo, geralmente exige uma ação para a esquerda.

\textbf{revisa aqui isaque.}

O procedimento automático atual utiliza sequências do giroscópio e do acelerômetro alinhadas aos timestamps do vídeo. A velocidade angular medida pelo giroscópio auxilia na identificação dos intervalos de manobra para a esquerda ou direita, enquanto os padrões do acelerômetro auxiliam na distinção entre caminhada contínua, desaceleração e parada. Os sinais são filtrados e agregados em uma janela temporal ao redor de cada quadro ou clipe curto. \todo{inserir filtro final, duração da janela, transformação dos eixos e limiares de decisão após a calibração. Tem que citar algum trabalho que fale dos metodos.} A anotação do protocolo permanece como referência no nível do evento, permitindo detectar e revisar divergências entre a ação planejada e os rótulos derivados do movimento.

Dada uma entrada visual $X_t$, representada por um quadro ou por um clipe curto encerrado no instante $t$, a tarefa direta é definida por

\begin{equation}
    f_{\theta}(X_t)=\hat{y}_t.
\end{equation}

Uma segunda formulação utiliza um extrator compartilhado com duas cabeças,

\begin{equation}
    f_{\theta}(X_t)=\left(\hat{y}_t,\hat{o}_t\right),
\end{equation}

que predizem conjuntamente a ação e a presença de obstáculo. A primeira configuração testa se o modelo pode aprender diretamente uma ação local segura. A segunda introduz a supervisão explícita de obstáculo como objetivo auxiliar e interpretável, seguindo a motivação multitarefa de DroNet e HDIN \cite{loquercio2018,chang2022}. Nesta etapa, \textsc{parar} representa a resposta de parada demonstrada no protocolo, e não a comprovação de que inexiste um caminho alternativo fora do campo de visão da câmera.

\subsection{Estatísticas da Base e Particionamento}

Após a validação das gravações, a versão atual da base reúne 592 eventos de navegação, distribuídos em 13 sessões e com duração total de 1~h, 45~min e 42~s. O IFCE concentra 320 eventos provenientes de seis sessões, a BECE contém 95 eventos de duas sessões e a UFC reúne 177 eventos de cinco sessões.

Na formulação de ação, todos os eventos possuem rótulo definido: 180 pertencem à classe \textsc{frente}, 197 à classe \textsc{esquerda}, 187 à classe \textsc{direita} e 28 à classe \textsc{parar}. Para a formulação binária de presença de obstáculo, 142 eventos representam caminho livre e 450 contêm algum obstáculo. Esses valores correspondem ao nível dos eventos completos de gravação; a quantidade de quadros ou clipes utilizados pelos modelos depende da política de amostragem temporal adotada. A Tabela~\ref{tab:stats} resume essa distribuição.

\begin{table}[t]
\caption{Distribuição dos rótulos no nível dos eventos válidos.}
\label{tab:stats}
\centering
\scriptsize
\begin{tabular}{lrr}
\toprule
Família & Classe & Eventos \\
\midrule
Ação & Frente & 180 \\
     & Esquerda & 197 \\
     & Direita & 187 \\
     & Parar & 28 \\
\midrule
Obstáculo & Livre & 142 \\
          & Obstáculo & 450 \\
\bottomrule
\end{tabular}
\end{table}

O particionamento é definido no nível das gravações, antes da extração de quadros ou da formação de clipes. Cada vídeo de origem e todas as amostras derivadas dele permanecem em um único subconjunto. Assim, modelos que realizam classificação quadro a quadro podem utilizar imagens individuais, enquanto modelos temporais podem receber sequências de quadros, sem que amostras provenientes da mesma gravação sejam distribuídas entre treino, validação e teste.

Além disso, gravações correlacionadas de um mesmo corredor ou configuração de cena são mantidas no mesmo subconjunto. Essa regra reduz a possibilidade de que o modelo seja avaliado em imagens quase idênticas às observadas durante o treinamento. A avaliação considera uma divisão sem sobreposição de cenas e um protocolo entre locais, no qual os modelos são treinados com dados de duas instituições e avaliados na terceira. \todo{inserir a composição final dos conjuntos de treino, validação e teste}

Essa estratégia é mais rigorosa do que a divisão aleatória por quadros. Em \cite{medeiros2021}, as imagens foram reorganizadas aleatoriamente e divididas em treinamento, validação e teste, sem que fosse explicitada uma separação por gravação ou cena. Em bases constituídas por quadros extraídos de vídeos, a divisão aleatória pode distribuir imagens com alta correlação espaço-temporal entre os subconjuntos, causando vazamento de informação e superestimando o desempenho dos modelos \cite{figueiredo2024leakage}. No \sensecv{}, todas as amostras derivadas de uma mesma gravação permanecem juntas, e a avaliação entre locais busca fornecer uma estimativa mais realista da capacidade de generalização para corredores, objetos e ambientes não observados durante o treinamento.

\subsection{Métricas de Avaliação}

O desempenho dos modelos é avaliado por meio da acurácia e do \emph{F1-score}. Essas métricas são calculadas a partir da matriz de confusão, considerando tanto a formulação multiclasse de ação quanto a formulação binária de presença de obstáculo. A acurácia representa a proporção total de amostras classificadas corretamente, enquanto o \emph{F1 score} combina precisão e revocação em uma única medida \cite{rainio2024metrics}.

Seja $C$ o número de classes e $\mathbf{M}\in\mathbb{N}^{C\times C}$ a matriz de confusão, em que $M_{ij}$ representa o número de amostras cuja classe real é $i$ e cuja classe prevista é $j$. A acurácia é definida por

\begin{equation}
\mathrm{Acc} =
\frac{\displaystyle\sum_{i=1}^{C} M_{ii}}
{\displaystyle\sum_{i=1}^{C}\sum_{j=1}^{C} M_{ij}}.
\label{eq:accuracy}
\end{equation}

Para cada classe $i$, a precisão e a revocação são calculadas,
respectivamente, por

\begin{equation}
\mathrm{Prec}_{i} =
\frac{M_{ii}}
{\displaystyle\sum_{k=1}^{C}M_{ki}},
\qquad
\mathrm{Rec}_{i} =
\frac{M_{ii}}
{\displaystyle\sum_{k=1}^{C}M_{ik}}.
\label{eq:precision_recall}
\end{equation}

O \emph{F1-score} da classe $i$ corresponde à média harmônica entre precisão e
revocação:

\begin{equation}
F_{1,i} =
2\frac{\mathrm{Prec}_{i}\mathrm{Rec}_{i}}
{\mathrm{Prec}_{i}+\mathrm{Rec}_{i}}.
\label{eq:f1_class}
\end{equation}

Como as classes da base apresentam frequências distintas, especialmente pela menor quantidade de eventos da classe \textsc{parar}, os resultados são reportados utilizando a macro-F1, definida como a média não ponderada dos valores obtidos para cada classe:

\begin{equation}
\mathrm{Macro}\text{-}F_{1} =
\frac{1}{C}\sum_{i=1}^{C}F_{1,i}.
\label{eq:macro_f1}
\end{equation}

Dessa forma, todas as classes contribuem igualmente para a métrica final, independentemente de sua frequência. A acurácia é apresentada como uma medida global complementar, enquanto a macro-F1 é utilizada como principal indicador para a comparação entre os modelos \cite{farhadpour2024metrics}.

\subsection{Configuração Experimental}

Foram avaliadas as arquiteturas MobileNetV1, EfficientNet-B0, ResNet-50 e
MobileNetV4 em uma formulação multitarefa de duas cabeças. Os extratores foram
inicializados com pesos pré-treinados no ImageNet, por meio das implementações
do pacote \textit{timm}: \texttt{mobilenetv1\_100.ra4\_e3600\_r224\_in1k},
\texttt{efficientnet\_b0.ra\_in1k}, \texttt{resnet50.a1\_in1k} e
\texttt{mobilenetv4\_conv\_small.e2400\_r224\_in1k}. A cada extrator foram
acopladas duas camadas lineares: uma cabeça binária para \textsc{livre} e
\textsc{obstáculo}, e uma cabeça de três classes para \textsc{esquerda},
\textsc{direita} e \textsc{nenhum desvio}. Os modelos desta comparação
realizam inferência quadro a quadro; os parâmetros de comprimento e passo de
sequência registrados nos scripts pertencem ao conjunto de dados exportado e
não são usados como entrada temporal por essas quatro arquiteturas.

Todas as imagens foram redimensionadas para $224\times224$ pixels por
interpolação bicúbica, convertidas em tensores e normalizadas com as médias
$(0{,}485, 0{,}456, 0{,}406)$ e os desvios padrão
$(0{,}229, 0{,}224, 0{,}225)$ do ImageNet. O treinamento utilizou lotes de
32 imagens, otimizador AdamW, taxa de aprendizado inicial de $10^{-4}$ e
decaimento de pesos de $10^{-4}$. A função objetivo foi a soma das perdas de
entropia cruzada das duas cabeças, calculadas somente nas amostras ativas para
a respectiva cabeça. A taxa de aprendizado foi reduzida por um fator de 0,5
após duas épocas sem redução da perda de validação, até o limite de $10^{-7}$.
O limite máximo foi de 40 épocas, com parada antecipada após seis épocas sem
melhora; o estado de menor perda de validação foi recuperado para a avaliação
final. O treinamento usou precisão mista quando uma GPU estava disponível e
semente fixa 43, com operações cuDNN determinísticas.

Para atenuar o desequilíbrio, as amostras de \textsc{caminho livre} na cabeça
de obstáculo e de \textsc{nenhum desvio} na cabeça lateral foram
subamostradas independentemente, de forma uniforme entre os clipes, até a
quantidade da respectiva classe minoritária. As amostras de treinamento
resultantes foram espelhadas horizontalmente; nessa transformação, os rótulos
\textsc{esquerda} e \textsc{direita} foram trocados, enquanto
\textsc{nenhum desvio} foi preservado. O mesmo balanceamento por cabeça foi
aplicado à validação, mas sem aumento de dados; o conjunto de teste permaneceu
inalterado. Assim, o balanceamento e o espelhamento não introduzem cópias
aumentadas na avaliação reportada.

Os relatórios de execução registram 9 épocas concluídas para MobileNetV1 e
MobileNetV4 e 7 para EfficientNet-B0 e ResNet-50, todas abaixo do limite
máximo devido ao critério de parada antecipada. As divisões exportadas foram
mantidas fixas entre os modelos, permitindo que as diferenças da
Tabela~\ref{tab:results_models} sejam atribuídas às arquiteturas sob a mesma
configuração de treinamento.

\section{Resultados}

Esta seção apresenta resultados preliminares da avaliação no conjunto de
teste, mantido separado por gravações, cenas e local de coleta. As métricas
são reportadas para as duas cabeças da formulação multitarefa: a cabeça de
\emph{obstáculo}, com as classes \textsc{livre} e \textsc{obstáculo}, e a
cabeça \emph{lateral}, com as classes \textsc{esquerda}, \textsc{direita} e
\textsc{nenhum desvio}. Em cada cabeça, o F1-score corresponde à macro-F1,
conforme definido na Seção III. Os valores da Tabela~\ref{tab:results_models}
foram calculados no conjunto de teste; portanto, não incluem amostras usadas
para treinamento, validação ou aumento de dados.

\begin{table}[t]
\caption{Resultados preliminares no conjunto de teste e número de parâmetros
treináveis das arquiteturas. Acc.: acurácia; F1: macro-F1.}
\label{tab:results_models}
\centering
\scriptsize
\renewcommand{\arraystretch}{1.12}
\setlength{\tabcolsep}{2pt}
\begin{tabular}{p{1.55cm}p{2.1cm}cc}
\toprule
\textbf{Modelo} &
\shortstack{\textbf{Parâmetros}\\\textbf{treináveis}} &
\shortstack{\textbf{Obstáculo (\%)}\\\textbf{Acc. / F1}} &
\shortstack{\textbf{Lateral (\%)}\\\textbf{Acc. / F1}} \\
\midrule
MobileNetV1 & 3\,212\,101 & 95,19 / 95,12 & \textbf{86,22 / 85,88} \\
EfficientNet-B0 & 4\,013\,953 & 93,90 / 93,87 & 70,73 / 67,28 \\
ResNet-50 & 23\,518\,277 & 89,45 / 89,35 & 60,60 / 53,60 \\
MobileNetV4 & 2\,497\,829 & \textbf{96,68 / 96,64} & 78,84 / 78,13 \\
\bottomrule
\end{tabular}
\end{table}

Os resultados indicam que a presença de obstáculo é a tarefa mais estável
entre as arquiteturas avaliadas. A MobileNetV4 apresentou o maior desempenho
nessa cabeça, alcançando 96,68\% de acurácia e 96,64\% de macro-F1, mesmo
sendo a rede com menor número de parâmetros. A MobileNetV1 também obteve
resultado elevado para obstáculo e foi a melhor arquitetura na predição
lateral, com 86,22\% de acurácia e 85,88\% de macro-F1. Assim, para esta
avaliação inicial, as arquiteturas móveis apresentaram a melhor relação entre
capacidade preditiva e tamanho do modelo.

Em contraste, o aumento do número de parâmetros não se traduziu em melhor
desempenho: a ResNet-50, apesar de possuir 23\,518\,277 parâmetros, obteve os
menores resultados nas duas cabeças. A diferença é mais pronunciada na cabeça
lateral, cuja macro-F1 variou de 53,60\% a 85,88\%. Esse comportamento sugere
que distinguir a direção do desvio a partir de quadros isolados é mais
sensível à variação entre locais, à aparência dos corredores e ao
desbalanceamento residual das classes do que a identificação binária de
obstáculos. Como se trata de resultados preliminares, essa interpretação deve
ser confirmada com repetições sob sementes distintas, análise por classe e
avaliação adicional dos erros nos diferentes ambientes.


\section{Conclusão}
Este artigo apresentou a primeira versão do \sensecv, uma base visual--inercial egocêntrica e multilocal para navegação assistiva local em ambientes internos. Um aplicativo Android dedicado registra vídeo, metadados e múltiplos dados de sensores do smartphone de forma sincronizada. O método atual de rotulagem utiliza acelerômetro e giroscópio, enquanto as demais modalidades são preservadas para pesquisas futuras. 

\begin{thebibliography}{99}

\bibitem{who2026}
World Health Organization,
``Blindness and vision impairment,''
Feb. 10, 2026. [Online]. Available:
\url{https://www.who.int/news-room/fact-sheets/detail/blindness-and-visual-impairment}.
[Accessed: Jul. 27, 2026].

\bibitem{casanova2025}
E. Casanova, D. Guffanti, and L. Hidalgo,
``Technological advancements in human navigation for the visually
impaired: A systematic review,''
\emph{Sensors}, vol. 25, no. 7, Art. no. 2213, 2025,
doi: 10.3390/s25072213.

\bibitem{khan2022}
D. Khan, Z. Cheng, H. Uchiyama, S. Ali, M. Asshad, and
K. Kiyokawa,
``Recent advances in vision-based indoor navigation: A systematic
literature review,''
\emph{Computers \& Graphics}, vol. 104, pp. 24--45, 2022,
doi: 10.1016/j.cag.2022.03.005.

\bibitem{gao2025}
Y. Gao \emph{et al.},
``A wearable obstacle avoidance device for visually impaired
individuals with cross-modal learning,''
\emph{Nature Communications}, vol. 16, Art. no. 2857, 2025,
doi: 10.1038/s41467-025-58085-x.

\bibitem{jadhav2025}
A. Jadhav \emph{et al.},
``AI Guide Dog: Egocentric path prediction on smartphone,''
\emph{Proceedings of the AAAI Symposium Series},
vol. 5, no. 1, pp. 220--227, 2025,
doi: 10.1609/aaaiss.v5i1.35591.

\bibitem{medeiros2021}
D. S. de Medeiros, T. H. Araújo, E. T. da Silva Júnior,
and G. L. B. Ramalho,
``Using images to avoid collisions and bypass obstacles in indoor
environments,''
in \emph{Anais Estendidos da XXXIV Conference on Graphics,
Patterns and Images}, 2021,
doi: 10.5753/sibgrapi.est.2021.20030.

\bibitem{figueiredo2024leakage}
R. B. D. Figueiredo and H. A. Mendes,
``Analyzing information leakage on video object detection datasets
by splitting images into clusters with high spatiotemporal
correlation,''
\emph{IEEE Access}, vol. 12, pp. 47646--47655, 2024,
doi: 10.1109/ACCESS.2024.3383047.

\bibitem{chitraranjan2025}
C. Chitraranjan, V. Vipulananthan, and T. Sritharan,
``Vision-based collision warning systems with deep learning:
A systematic review,''
\emph{Journal of Imaging}, vol. 11, no. 2, Art. no. 64, 2025,
doi: 10.3390/jimaging11020064.

\bibitem{breve2020}
F. A. Breve and C. N. Fischer,
``Visually impaired aid using convolutional neural networks,
transfer learning, and particle competition and cooperation,''
in \emph{Proc. 2020 International Joint Conference on Neural
Networks (IJCNN)}, 2020, pp. 1--8,
doi: 10.1109/IJCNN48605.2020.9207606.

\bibitem{bastidas2025detectra}
G. K. Bastidas-Guacho, M. A. Paguay Alvarado,
P. X. Moreno-Vallejo, P. R. Moreno-Costales,
N. S. Ocaña Yanza, and J. C. Troya Cuestas,
``Computer vision-based obstacle detection mobile system for
visually impaired individuals,''
\emph{Multimodal Technologies and Interaction},
vol. 9, no. 5, Art. no. 48, 2025,
doi: 10.3390/mti9050048.

\bibitem{plikynas2022}
D. Plikynas, A. Indriulionis, A. Laukaitis, and
L. Sakalauskas,
``Indoor-guided navigation for people who are blind:
Crowdsourcing for route mapping and assistance,''
\emph{Applied Sciences}, vol. 12, no. 1, Art. no. 523, 2022,
doi: 10.3390/app12010523.

\bibitem{loquercio2018}
A. Loquercio, A. I. Maqueda, C. R. del-Blanco, and
D. Scaramuzza,
``DroNet: Learning to fly by driving,''
\emph{IEEE Robotics and Automation Letters},
vol. 3, no. 2, pp. 1088--1095, 2018,
doi: 10.1109/LRA.2018.2795643.

\bibitem{chang2022}
Y. Chang, Y. Cheng, J. Murray, S. Huang, and G. Shi,
``The HDIN dataset: A real-world indoor UAV dataset with
multi-task labels for visual-based navigation,''
\emph{Drones}, vol. 6, no. 8, Art. no. 202, 2022,
doi: 10.3390/drones6080202.

\bibitem{samadzadeh2024}
A. Samadzadeh, M. H. Mojab, H. Soudani,
S. H. Mireshghollah, and A. Nickabadi,
``Amirkabir campus dataset: Real-world challenges and scenarios
of visual inertial odometry for visually impaired people,''
\emph{arXiv preprint arXiv:2401.03604}, 2024,
doi: 10.48550/arXiv.2401.03604.

\bibitem{karnan2022scand}
H. Karnan \emph{et al.},
``Socially CompliAnt Navigation Dataset (SCAND):
A large-scale dataset of demonstrations for social navigation,''
\emph{IEEE Robotics and Automation Letters},
vol. 7, no. 4, pp. 11807--11814, 2022,
doi: 10.1109/LRA.2022.3184025.

\bibitem{liu2019travel}
X. Liu, S. Zhang, J. Zeng, and F. Fan,
``Analysis and optimization strategy of travel system for urban
visually impaired people,''
\emph{Sustainability}, vol. 11, no. 6,
Art. no. 1735, 2019,
doi: 10.3390/su11061735.

\bibitem{yan2020flobot}
Z. Yan, S. Schreiberhuber, G. Halmetschlager,
T. Duckett, M. Vincze, and N. Bellotto,
``Robot perception of static and dynamic objects with an
autonomous floor scrubber,''
\emph{Intelligent Service Robotics},
vol. 13, no. 3, pp. 403--417, 2020,
doi: 10.1007/s11370-020-00324-9.

\bibitem{rainio2024metrics}
O. Rainio, J. Teuho, and R. Klén,
``Evaluation metrics and statistical tests for machine learning,''
\emph{Scientific Reports}, vol. 14,
Art. no. 6086, 2024,
doi: 10.1038/s41598-024-56706-x.

\bibitem{farhadpour2024metrics}
S. Farhadpour, T. A. Warner, and A. E. Maxwell,
``Selecting and interpreting multiclass loss and accuracy
assessment metrics for classifications with class imbalance:
Guidance and best practices,''
\emph{Remote Sensing}, vol. 16, no. 3,
Art. no. 533, 2024,
doi: 10.3390/rs16030533.

\end{thebibliography}

\end{document}
